\documentclass[11pt,a4paper]{article}

\usepackage[T1]{fontenc}
\usepackage[utf8]{inputenc}
\usepackage[margin=24mm]{geometry}
\usepackage{amsmath,amssymb,bm}
\usepackage{booktabs,longtable,array,tabularx}
\usepackage{graphicx}
\usepackage{xcolor}
\usepackage{listings}
\usepackage{hyperref}

\definecolor{codeblue}{RGB}{32,74,135}
\definecolor{codegreen}{RGB}{35,110,55}
\definecolor{codegray}{RGB}{90,90,90}
\definecolor{shade}{RGB}{247,247,247}

\hypersetup{
  colorlinks=true,
  linkcolor=codeblue,
  urlcolor=codeblue,
  citecolor=codeblue,
  pdftitle={Structural Remeshing and Modal Analysis of the Chang Wing--Propeller Model}
}

\lstdefinelanguage{Julia}{
  morekeywords={function,end,for,in,if,else,elseif,return,begin,let,struct,
    mutable,where,using,import,const,true,false,zeros,ones,eachindex},
  sensitive=true,
  morecomment=[l]{\#},
  morestring=[b]"
}

\lstset{
  language=Julia,
  basicstyle=\ttfamily\small,
  keywordstyle=\color{codeblue}\bfseries,
  commentstyle=\color{codegreen},
  stringstyle=\color{codegray},
  backgroundcolor=\color{shade},
  frame=single,
  rulecolor=\color{black!20},
  columns=fullflexible,
  keepspaces=true,
  breaklines=true,
  showstringspaces=false,
  tabsize=4
}

\newcommand{\mat}[1]{\bm{#1}}
\newcommand{\vect}[1]{\bm{#1}}
\newcommand{\Hz}{\,\mathrm{Hz}}

\title{Structural Remeshing and Modal Analysis\\
of the Chang Wing--Propeller Model}
\author{Technical verification note}
\date{2 September 2026}

\begin{document}

\maketitle

\begin{abstract}
This note documents the structural interpretation and remeshing procedure for
the Chang wing--propeller model.  Particular attention is given to the
different coordinate conventions, the distinction between distributed
stiffness and lumped inertia, the parallel-axis contribution in the nodal
spatial-inertia matrix, and the source of the failure observed with a
20-element uniform mesh.  Modal results are presented for the isolated wing
and for the coupled wing--rigid-pylon--propeller system.  The central result is
that a 30-element structural mesh is adequate when stiffness is integrated as
a distributed field and the complete spatial inertia is conservatively
lumped at target nodal control volumes.  The first five isolated-wing modes
and the first ten coupled modes are then within approximately \(0.30\%\) of a
320-element numerical check.
\end{abstract}

\tableofcontents
\newpage

\section{Purpose and scope}

The original model-property workbook uses \(X\) in the spanwise direction,
\(Y\) in the chordwise direction, and \(Z\) in the vertical direction.  The
present structural implementation uses the nodal degree-of-freedom order
\begin{equation}
\vect{q}_n =
\begin{bmatrix}
u_s & v_c & w_z & \theta_s & \theta_c & \theta_z
\end{bmatrix}^{\mathsf T},
\end{equation}
where the subscripts denote spanwise, chordwise, and vertical quantities.
The inertia components in the workbook must therefore first be retained in
their literal physical axes; any structural-to-aerodynamic basis change must
be applied only once at the coupling interface.

The objectives of this note are to:
\begin{enumerate}
  \item distinguish the distributed stiffness data from the concentrated
        mass and inertia data;
  \item define a physically consistent nodal rigid-body mass matrix;
  \item explain the limitations of direct componentwise interpolation;
  \item identify why an intermediate remap converged very slowly;
  \item quantify structural mesh convergence for the isolated and coupled
        systems; and
  \item recommend a practical discretization for the time-domain UVLM model.
\end{enumerate}

The numerical source is \texttt{Model Properties.xlsx}, in particular the
worksheets \texttt{Wing Construction} and \texttt{Wing stiffness}.  The
standalone numerical implementation used for the audit is
\path{experiments/conservative_lumped_modal/ConservativeLumpedModal.jl}.
The same conservative formulation is now also active in the production Chang
case files.

\section{Interpretation of the workbook data}

\subsection{Flexible beam and rigid inertia elements}

The connectivity in \texttt{Wing Construction} is essential.  Flexible
elements 20--35 form the spanwise beam.  In contrast, elements 52--67 connect
each non-root flexible-wing station to an offset node at the same span
coordinate.  These are rigid mass-attachment elements.  Consequently, the
16 tabulated mass, center-of-gravity, and inertia records are concentrated
nodal rigid-body properties in the reference finite-element model; they are
not 16 distributed flexible-element densities.

There are 17 flexible-wing stations including the fixed root.  The root has
no independent inertia record, while the other 16 stations carry the
tabulated inertias.  Their total mass after unit conversion is
\begin{equation}
M_{\mathrm{wing}} = 177.32851668\ \mathrm{kg}.
\end{equation}
The corresponding spanwise centroid of the 16 discrete source lumps is
\begin{equation}
\bar{x}_{\mathrm{source}} = 3.549179\ \mathrm{m}.
\end{equation}

The source grid is nonuniform.  A uniform 16-element mesh is therefore not
the same discretization as the 16-element workbook mesh, even though the two
meshes have the same number of elements.  This distinction is important when
interpreting the modal tables.

\subsection{Distributed stiffness}

The stiffness entries define a spanwise constitutive field.  For the strain
ordering used by the beam element, the sectional matrix is
\begin{equation}
\mat{C}(x)=
\begin{bmatrix}
EA & 0 & 0 & 0\\
0 & EI_z & EI_{zy} & 0\\
0 & EI_{zy} & EI_y & 0\\
0 & 0 & 0 & GJ
\end{bmatrix}.
\end{equation}
The \(EI_{zy}\) term is retained and provides bending cross-coupling.

Several stiffness stations occur in near-equal pairs, for example
\(\eta=0.04443\) and \(0.04444\).  These pairs represent sharp property
changes rather than a broad smooth transition.  This makes a single
element-center sample sensitive to the alignment between a uniform target
mesh and the stiffness breakpoints.

\section{Structural stiffness assembly}

\subsection{Element-center approximation}

The current production mapping evaluates the interpolated stiffness at the
center of each target element,
\begin{equation}
\mat{C}_e \approx \mat{C}(x_{e,c}),
\end{equation}
and then treats it as constant over the element.  This is equivalent to
approximating
\begin{equation}
\mat{K}_e = \int_{x_e}^{x_{e+1}}
\mat{B}(x)^{\mathsf T}\mat{C}(x)\mat{B}(x)\,dx
\end{equation}
with a property value taken at one point.

This approximation is appropriate when an element lies inside a constant or
slowly varying stiffness region.  It is less reliable when an element spans a
sharp property boundary.  As the number of uniform elements changes, a
boundary can move from one side of an element center to the other.  The modal
sequence can therefore oscillate with mesh count rather than converge
monotonically.

\subsection{Breakpoint-aware stiffness integration}

The mesh-independent treatment used in this audit is:
\begin{enumerate}
  \item determine every workbook stiffness breakpoint lying inside the
        target element;
  \item split the element integration interval at those breakpoints;
  \item evaluate the interpolated complete matrix \(\mat C(x)\) at Gauss
        points within each smooth subinterval; and
  \item accumulate \(\mat B^{\mathsf T}\mat C\mat B\) using the physical
        subinterval Jacobian.
\end{enumerate}

Thus the structural mesh need not be aligned with the workbook stiffness
stations.  At 30 elements, center sampling happens to agree almost exactly
with breakpoint-aware integration.  At 20 and 32 elements, center sampling
can change individual modes by approximately \(1\%\), showing that the
apparent oscillation is a quadrature/alignment effect rather than a need for
hundreds of beam elements.

\section{Lumped inertia and spatial-inertia consistency}

\subsection{Source tributary lengths}

Let the source nodes be \(x_i\).  The nodal tributary lengths used in the
original interpolation are
\begin{align}
L_1 &= \frac{x_2-x_1}{2},\\
L_i &= \frac{(x_i-x_{i-1})+(x_{i+1}-x_i)}{2},
       \qquad 2\le i\le N,\\
L_{N+1} &= \frac{x_{N+1}-x_N}{2}.
\end{align}
For an integrated nodal quantity \(a_i\), the reconstructed nodal density is
\begin{equation}
\rho_{a,i}=\frac{a_i}{L_i}.
\end{equation}

The original code interpolated the density at each new node and multiplied it
by the target tributary length:
\begin{equation}
a_j^{\mathrm{new}} =
L_j^{\mathrm{new}}\,
\mathcal{I}\!\left[\rho_a\right](x_j^{\mathrm{new}}),
\end{equation}
where \(\mathcal I\) denotes one-dimensional linear interpolation.

\begin{lstlisting}[caption={Original density-interpolation pattern.}]
rho_M_orig = M_orig_nodes ./ L_node_orig
rho_M_new  = interpolate_1d(X_nodes_orig, rho_M_orig, X_nodes_new)
m_node_vec = rho_M_new .* L_node_new
\end{lstlisting}

This is a reasonable lumped quadrature and exactly reproduces each source
value when the target nodes equal the source nodes.  It is not, however,
exactly conservative on a different grid.  At 30 uniform elements it gives
\(175.41043878\ \mathrm{kg}\), which is \(1.91807790\ \mathrm{kg}\), or
\(1.082\%\), below the source mass.

\subsection{Why mass, CG, and inertia cannot be remapped independently}

Let \(m\) be a nodal mass, \(\vect\eta\) the vector from the beam reference
axis to the mass center, and \(\mat J_{CG}\) the inertia tensor about the
center of mass.  The first moment is
\begin{equation}
\vect h = m\vect\eta.
\end{equation}
If \(m\), \(\vect\eta\), and \(\mat J_{CG}\) are interpolated separately, the
interpolated product generally does not satisfy
\begin{equation}
\mathcal I[m\vect\eta]
= \mathcal I[m]\,\mathcal I[\vect\eta].
\end{equation}
The same problem appears more strongly in the quadratic parallel-axis terms.
Therefore, independently interpolated CG and inertia arrays do not generally
describe the same integrated rigid body.

The correct conserved quantities are the mass, the three first moments, and
the complete inertia about a common reference point.  These are naturally
combined in one spatial-inertia matrix.

\subsection{Complete \(6\times6\) spatial inertia}

Define the skew-symmetric cross-product matrix
\begin{equation}
\widetilde{\vect\eta}=
\begin{bmatrix}
0 & -\eta_z & \eta_y\\
\eta_z & 0 & -\eta_x\\
-\eta_y & \eta_x & 0
\end{bmatrix}.
\end{equation}
For nodal generalized velocity

\begin{equation}
\vect v_n =
\begin{bmatrix}
\dot{\vect u}_n\\ \dot{\vect\theta}_n
\end{bmatrix},
\end{equation}
the rigid-body kinetic energy is
\begin{equation}
T_n=\frac{1}{2}\vect v_n^{\mathsf T}\mat H_n\vect v_n,
\end{equation}
where
\begin{equation}
\boxed{
\mat H_n=
\begin{bmatrix}
m\mat I_3 & -m\widetilde{\vect\eta}\\
m\widetilde{\vect\eta} &
\mat J_{CG}-m\widetilde{\vect\eta}^{2}
\end{bmatrix}.}
\label{eq:spatial-inertia}
\end{equation}
Because
\begin{equation}
-m\widetilde{\vect\eta}^{2}
=m\left[(\vect\eta^{\mathsf T}\vect\eta)\mat I_3
-\vect\eta\vect\eta^{\mathsf T}\right],
\end{equation}
the lower-right block in Equation~\eqref{eq:spatial-inertia} is exactly the
parallel-axis shift from the CG to the beam node.

The Schur complement of the translational block \(m\mat I_3\) is
\(\mat J_{CG}\).  Thus, for \(m>0\), a physically positive-definite
\(\mat J_{CG}\) produces a physically valid spatial-inertia block.

All products of inertia, \(I_{xy}\), \(I_{xz}\), and \(I_{yz}\), are retained
in \(\mat J_{CG}\).  No diagonal-only approximation is made.

\subsection{Correction to the displayed Julia block}

If the supplied inertia is about the CG, the following lower-right block is
incomplete:

\begin{lstlisting}
ms_W = [m_val * I3              -m_val * chang_tilde(eta);
        m_val * chang_tilde(eta)  inertia_matrix]
\end{lstlisting}

The consistent implementation is:

\begin{lstlisting}[caption={Nodal spatial inertia when the input inertia is about the CG.}]
S = chang_tilde(eta)
J_at_node = inertia_matrix - m_val * S * S

ms_W = [m_val * I3   -m_val * S;
        m_val * S     J_at_node]
\end{lstlisting}

The current \texttt{chang\_point\_mass\_matrix} helper applies this
parallel-axis term when
\texttt{inertia\_reference = :center\_of\_mass}.

The modal consequence is substantial.  On the source mesh, treating the
spreadsheet inertia as CG inertia and applying the parallel-axis term gives a
first torsional mode of \(42.87465\Hz\).  Inserting the same CG inertia
directly into the lower-right block gives \(50.14898\Hz\).  The latter does
not reproduce the reference torsional behavior.

\section{Comparison of remapping strategies}

\subsection{Point-sampled density interpolation}

The original method samples each reconstructed density at the target nodes
and applies target tributary weights.  Its advantages are simplicity, a
positive mass matrix for the tested meshes, and exact reproduction when the
source grid itself is requested.  Its disadvantages are quadrature error and
the lack of exact conservation on an arbitrary target mesh.

The method becomes more consistent if complete spatial blocks, rather than
CG and inertia components independently, are interpolated.  Even then,
point-sampled quadrature at 30 elements retains the \(1.082\%\) total-inertia
error reported above.

\subsection{Independent signed-total rescaling}

An attempted correction rescales each interpolated scalar field independently:
\begin{equation}
a_j^{\mathrm{scaled}} = a_j^{\mathrm{raw}}
\frac{\sum_i a_i^{\mathrm{source}}}
     {\sum_j a_j^{\mathrm{raw}}}.
\end{equation}
This is unsafe for signed first moments and products of inertia.  Their sums
can be small because positive and negative contributions cancel.  A modest
quadrature error in the denominator can therefore create a very large scale
factor.

At 20 elements the observed consequences are:
\begin{itemize}
  \item maximum recovered CG offset: \(4.146195\ \mathrm{m}\);
  \item minimum eigenvalue of a recovered CG inertia:
        \(-55.46860\ \mathrm{kg\,m^2}\); and
  \item minimum eigenvalue of the global free mass matrix:
        \(-18.74893\).
\end{itemize}
The generalized eigenvalue solve consequently raises a
\texttt{PosDefException}.  This failure occurs before Generalized-
\(\alpha\) time integration and is not caused by the GA parameters or
coupling relaxation.

\subsection{Element-integrated density lumped at the outer node}

An intermediate conservative method evaluated
\begin{equation}
\mat H_{j+1}=\int_{x_j}^{x_{j+1}}\mat\rho_H(x)\,dx
\end{equation}
and placed the entire integral at the outer node \(x_{j+1}\).  Although it
conserves the sum of the spatial-inertia blocks and preserves positive
definiteness, it shifts each element's inertia approximately one-half element
toward the tip.

The effect can be seen in the mass centroid:
\begin{center}
\begin{tabular}{lcc}
\toprule
Representation & Total mass (kg) & Spanwise centroid (m)\\
\midrule
Source 16 discrete lumps & 177.32852 & 3.54918\\
Outer-node remap, 16 elements & 177.32852 & 3.78647\\
Outer-node remap, 20 elements & 177.32852 & 3.73939\\
Outer-node remap, 30 elements & 177.32852 & 3.67689\\
Control-volume remap, 30 elements & 177.32852 & 3.55564\\
\bottomrule
\end{tabular}
\end{center}

This explains the slowly increasing frequency sequence
\(6.22470, 6.30302, 6.35567, 6.40404,\ldots\).  That sequence is not produced
by the original point-sampled method.  It is a first-order mass-location bias
created by outer-node lumping.

\subsection{Conservative control-volume spatial remap}

The corrected experimental method first constructs the complete source
spatial blocks and divides them by their source tributary lengths:
\begin{equation}
\mat\rho_{H,i}=\frac{\mat H_i}{L_i}.
\end{equation}
The complete matrix density is linearly interpolated between source stations.
For target node \(j\), define its control-volume boundaries by adjacent
element midpoints,
\begin{equation}
x_{j,-}=\frac{x_{j-1}+x_j}{2}, \qquad
x_{j,+}=\frac{x_j+x_{j+1}}{2},
\end{equation}
with the physical root and tip used as the end boundaries.  The target lump
is
\begin{equation}
\mat H_j^{\mathrm{new}} =
\int_{x_{j,-}}^{x_{j,+}}\mat\rho_H(x)\,dx.
\end{equation}
The integral is evaluated exactly over each overlap with a source interpolation
interval.  All weights are nonnegative, so a nonnegative combination of
valid source blocks remains valid.  The complete spatial-inertia sum is
conserved to machine precision.

The small control-volume contribution associated with the fixed root is
transferred to the first free node.  This retains the complete reconstructed
mass in the free structural model and avoids a massless free node.  The fixed
root itself has no dynamic degree of freedom.

\subsection{An unavoidable modeling distinction}

The direct \texttt{source\_16} model is a discrete set of point inertias.  A
uniform refined model in which every free node has a positive lumped inertia
is necessarily a distributed approximation.  These requirements cannot all
be satisfied simultaneously without choosing a continuous reconstruction of
the source data:
\begin{itemize}
  \item retaining only the exact 16 point inertias preserves the original
        discrete model, but added target nodes can be massless;
  \item spreading the inertias gives every free node a valid mass block, but
        defines a continuous mass distribution that is not identical to the
        original point-mass model at high spatial frequencies.
\end{itemize}

For this reason, the refined uniform-mesh frequencies should converge to the
chosen reconstructed distribution, not exactly to every high-frequency mode
of the coarse 16-lump model.  The low modes remain very close because their
mode shapes vary slowly over the source station spacing.

For transparency, if the control-volume remap is applied on the original
nonuniform source grid instead of using the source blocks directly, the first
frequency is \(6.61462\Hz\), versus \(6.62627\Hz\) for the direct source
lumps.  The report therefore labels the direct discrete model and the
continuous uniform remesh as distinct representations.

\section{Modal analysis procedure}

\subsection{Isolated wing}

After assembly and elimination of the six fixed-root degrees of freedom, the
undamped modes solve
\begin{equation}
\mat K\vect\phi_r = \omega_r^2\mat M\vect\phi_r,
\qquad
f_r=\frac{\omega_r}{2\pi}.
\end{equation}
Both \(\mat M\) and \(\mat K\) are symmetrized before the generalized
eigenvalue solution.  The mass matrix must be positive definite after root
elimination.  Modes are classified from the relative modal content in the
out-of-plane, in-plane, torsional, and axial degree-of-freedom groups.

\subsection{Coupled wing--pylon--propeller system}

The structural comparison adds two rigid-pylon generalized coordinates,
pitch and yaw.  Their uncoupled nominal frequency is \(7.97\Hz\), with
\begin{equation}
\mat K_p =
\begin{bmatrix}
19220 & 0\\ 0 & 18916
\end{bmatrix},
\qquad
\mat M_p =
\begin{bmatrix}
I_{n\theta} & 0\\ 0 & I_{n\psi}
\end{bmatrix},
\end{equation}
where \(I_{n\theta}\) and \(I_{n\psi}\) are selected to produce the nominal
pitch and yaw frequencies.  The pylon and propeller translational mass,
first moments, root rotational inertias, and pitch/yaw cross-inertias are
included in the coupled mass matrix.

The attachment is maintained at the exact span fraction
\(\eta_p=0.83\) by work-conjugate interpolation between adjacent structural
nodes.  This removes nearest-node snapping from the structural mesh study.
The modal table is undamped and non-gyroscopic; aerodynamic and gyroscopic
effects belong to the subsequent time-domain analysis.

\section{Numerical results}

\subsection{Thirty-element remap comparison}

\begin{table}[htbp]
\centering
\caption{Comparison of three interpretations on a 30-element uniform mesh.}
\resizebox{\textwidth}{!}{%
\begin{tabular}{lrrrrrrrr}
\toprule
Method & Mass (kg) & Spatial error & min eig. \(M\) & OOP1 & IP1 & OOP2 & T1 & OOP3\\
\midrule
Point-sampled complete density
 & 175.41044 & \(1.082\times10^{-2}\) & 0.042964
 & 6.63808 & 20.81901 & 36.45316 & 42.89074 & 93.40783\\
Control-volume complete density
 & 177.32852 & \(1.60\times10^{-16}\) & 0.049233
 & 6.63049 & 20.79605 & 36.47010 & 42.84774 & 93.52212\\
Preceding-element interpretation
 & 177.32852 & \(8.01\times10^{-17}\) & 0.042964
 & 6.97171 & 21.85155 & 37.48588 & 44.24148 & 94.78981\\
\bottomrule
\end{tabular}}
\end{table}

The preceding-element interpretation treats each inertia record as the
integral of the preceding flexible spanwise element.  It is included only as
a sensitivity calculation and conflicts with the rigid-element connectivity
in the workbook.

\subsection{Isolated-wing natural frequencies}

\begin{table}[htbp]
\centering
\caption{Selected isolated-wing natural frequencies in Hz.}
\begin{tabular}{lrrrrr}
\toprule
Case & OOP1 & IP1 & OOP2 & T1 & OOP3\\
\midrule
Source 16, direct lumps & 6.62627 & 20.78295 & 36.38983 & 42.87465 & 92.42593\\
Uniform 16 & 6.63403 & 20.80353 & 36.40711 & 42.87282 & 92.79187\\
Uniform 20 & 6.63615 & 20.81031 & 36.46820 & 42.90134 & 93.25806\\
Uniform 24 & 6.63614 & 20.81087 & 36.48578 & 42.89387 & 93.40853\\
Uniform 30 & 6.63049 & 20.79605 & 36.47010 & 42.84774 & 93.52212\\
Uniform 40 & 6.63621 & 20.81251 & 36.52306 & 42.90235 & 93.69693\\
Uniform 80 & 6.63582 & 20.81170 & 36.52569 & 42.88710 & 93.78645\\
Uniform 160 & 6.63476 & 20.80894 & 36.52376 & 42.87981 & 93.80210\\
Uniform 320 & 6.63440 & 20.80793 & 36.52059 & 42.87411 & 93.80168\\
Published reference & 6.63000 & 20.93000 & 36.66000 & 43.50000 & 93.31000\\
\bottomrule
\end{tabular}
\end{table}

For the 30-element model, the differences from the published values are
\begin{center}
\begin{tabular}{lrrrrr}
\toprule
Mode & OOP1 & IP1 & OOP2 & T1 & OOP3\\
\midrule
Difference (\%) & \(+0.007\) & \(-0.640\) & \(-0.518\) & \(-1.499\) & \(+0.227\)\\
\bottomrule
\end{tabular}
\end{center}
The torsional difference does not disappear with structural refinement: the
320-element result remains \(1.439\%\) below the published value.  It is
therefore associated with the adopted model data or modeling assumptions,
not insufficient structural mesh resolution.

\subsection{Coupled natural frequencies}

\begin{table}[htbp]
\centering
\caption{First ten undamped coupled frequencies in Hz.}
\resizebox{\textwidth}{!}{%
\begin{tabular}{lrrrrrrrrrr}
\toprule
Case & $f_1$ & $f_2$ & $f_3$ & $f_4$ & $f_5$ & $f_6$ & $f_7$ & $f_8$ & $f_9$ & $f_{10}$\\
\midrule
Source 16
& 5.90034 & 7.95783 & 8.52744 & 19.38857 & 35.10018
& 39.27353 & 91.01576 & 108.18211 & 114.09417 & 155.65989\\
Uniform 30
& 5.90287 & 7.95783 & 8.52882 & 19.39934 & 35.15803
& 39.27965 & 92.09443 & 108.36857 & 114.52226 & 156.74651\\
Uniform 320
& 5.90521 & 7.95779 & 8.52978 & 19.40920 & 35.18707
& 39.31663 & 92.37101 & 108.49956 & 114.60404 & 156.85025\\
\bottomrule
\end{tabular}}
\end{table}

The \(7.96\Hz\) pair is dominated by the rigid-pylon pitch/yaw coordinates.
The frequencies near the isolated-wing bending modes are shifted by the
pylon and propeller generalized mass and their coupling with the wing.

\subsection{Structural mesh convergence}

Table~\ref{tab:convergence} reports, for each uniform mesh, the largest
absolute relative difference among the first five isolated-wing frequencies
and among the first ten coupled frequencies, using the 320-element model as a
numerical check.

\begin{table}[htbp]
\centering
\caption{Maximum modal differences relative to the 320-element uniform mesh.}
\label{tab:convergence}
\begin{tabular}{rrr}
\toprule
Elements & Wing modes (\%) & Coupled modes (\%)\\
\midrule
16  & 1.0765 & 0.9689\\
20  & 0.5795 & 0.4805\\
24  & 0.4191 & 0.4142\\
30  & 0.2980 & 0.2994\\
32  & 0.1976 & 0.1386\\
36  & 0.1779 & 0.1670\\
40  & 0.1117 & 0.0894\\
48  & 0.0837 & 0.0660\\
60  & 0.0752 & 0.0685\\
80  & 0.0303 & 0.0409\\
120 & 0.0226 & 0.0182\\
160 & 0.0133 & 0.0174\\
240 & 0.0127 & 0.0157\\
320 & 0.0000 & 0.0000\\
\bottomrule
\end{tabular}
\end{table}

The convergence is not perfectly monotonic because the source distribution
contains sharp changes and the target grids do not all place nodes at the
same locations.  Nevertheless, the error envelope decreases rapidly.  A
30-element mesh is already below \(0.30\%\) for both comparisons.

\section{Recommended implementation}

For the current UVLM aeroelastic calculations, the recommended structural
procedure is:
\begin{enumerate}
  \item retain the workbook inertia axes and signs in their original physical
        coordinate system;
  \item convert every source mass, CG, and CG inertia tensor into the complete
        spatial block in Equation~\eqref{eq:spatial-inertia};
  \item divide complete blocks by the source tributary lengths and interpolate
        the matrix density, rather than independently scaling signed scalar
        components;
  \item integrate the density over target nodal control volumes and lump it at
        those nodes;
  \item verify total mass and the sum of the complete spatial-inertia matrix;
  \item require every free nodal block and the assembled free mass matrix to
        be positive definite;
  \item integrate \(\mat B^{\mathsf T}\mat C(x)\mat B\) across every target
        element, splitting at all stiffness breakpoints; and
  \item attach the pylon at its exact physical span location using
        work-conjugate interpolation, rather than snapping it to the nearest
        structural node.
\end{enumerate}

The following checks should be executed before any time-domain analysis:
\begin{align}
\epsilon_H &=
\frac{\left\|\sum_j\mat H_j^{\mathrm{new}}
-\sum_i\mat H_i^{\mathrm{source}}\right\|_\infty}
{\left\|\sum_i\mat H_i^{\mathrm{source}}\right\|_\infty},\\
\lambda_{\min}(\mat H_j) &> 0
\quad\text{for every free node},\\
\lambda_{\min}(\mat M_{\mathrm{free}}) &> 0.
\end{align}
For the 30-element conservative model,
\(\epsilon_H=1.60\times10^{-16}\) and
\(\lambda_{\min}(\mat M_{\mathrm{free}})=0.049233>0\).

\subsection{Production-code integration}

The default spanwise discretization is set to 30 elements in
\path{chang_case.jl}.  The production implementation is divided as follows:
\begin{itemize}
  \item \path{chang_model_parameters.jl} constructs and conservatively remaps
        the complete nodal spatial-inertia blocks;
  \item \path{chang_structural_matrices.jl} performs breakpoint-aware
        stiffness integration and applies the exact attachment operator;
  \item \path{chang_structural_model.jl} selects these corrected mass and
        stiffness paths for the active center-of-mass interpretation;
  \item \path{chang_uvlm_coupling.jl} uses the same interpolation weights for
        propeller motion and for the transpose load map; and
  \item \path{run_chang_linear_aeroelastic.jl} checks positive definiteness
        and matrix symmetry before beginning UVLM time marching.
\end{itemize}

At \(\eta_p=0.83\), the 30-element mesh brackets the attachment with nodes 25
and 26 and weights 0.1 and 0.9.  The resulting physical position is exactly
\(6.225\ \mathrm{m}\), independent of nearest-node rounding.

\section{Conclusions}

\begin{enumerate}
  \item The original idea of dividing nodal integrated properties by source
        tributary lengths and reconstructing a spanwise density is physically
        reasonable.
  \item The old frequency sequence beginning at \(6.22470\Hz\) was caused by
        lumping each target-element integral at its outer node.  It is not a
        valid indication that 160 elements are required.
  \item Mass, first moments, and rotational inertia must be remapped as one
        complete spatial-inertia object.  Independent scaling of signed
        components caused the indefinite 20-element mass matrix.
  \item The parallel-axis contribution is mandatory when the tabulated
        inertia is about the CG.  Omitting it raises the calculated first
        torsional frequency from approximately \(42.87\Hz\) to
        \(50.15\Hz\).
  \item Breakpoint-aware stiffness integration avoids mesh-alignment
        oscillations caused by the sharply stepped stiffness distribution.
  \item With the corrected formulation, 30 structural elements provide a
        practical and accurate mesh.  Using 160 elements is unnecessary for
        the present structural modal accuracy target.
  \item Exact preservation of the coarse 16-point-mass model and positive
        inertia at every node of an arbitrarily refined mesh are different
        modeling objectives.  The continuous reconstruction must be stated
        explicitly when reporting refined-mesh results.
\end{enumerate}

\appendix

\section{Reproduction files}

The numerical results in this note can be regenerated from the repository
root with:

\begin{lstlisting}[language={},caption={PowerShell command for the modal audit.}]
julia --project=lib/WingPropellerUVLM \
  lib/WingPropellerUVLM/examples/chang_linear_aeroelastic/\
  experiments/conservative_lumped_modal/run_modal_comparison.jl
\end{lstlisting}

The principal files are:
\begin{itemize}
  \item \path{experiments/conservative_lumped_modal/ConservativeLumpedModal.jl};
  \item \path{experiments/conservative_lumped_modal/run_modal_comparison.jl};
  \item \path{experiments/conservative_lumped_modal/output/modal_frequency_comparison.csv};
  \item \path{experiments/conservative_lumped_modal/output/structural_mesh_convergence.csv};
  \item \path{experiments/conservative_lumped_modal/output/mass_matrix_audit.csv}; and
  \item \path{experiments/conservative_lumped_modal/output/remap_method_comparison_30.csv}.
\end{itemize}

\end{document}
